\section{Quá trình Sử dụng Trợ lý AI}

Trong quá trình làm đồ án, nhóm dùng trợ lý lập trình \textbf{Claude Code}
(mô hình Claude của Anthropic) như một cộng sự để tăng tốc xây dựng hệ
thống và báo cáo. Phần này không lặp lại danh sách mô hình AI bên trong
sản phẩm (đã trình bày ở Chương~7 và Chương~9), mà \textbf{trình bày cách
nhóm đã prompt và với mục đích gì}. Nội dung dưới đây được tổng hợp từ
nhật ký prompt mà nhóm lưu lại qua các phiên làm việc.

\subsection{Các Nhóm Prompt và Mục đích}
\contrib

Việc sử dụng AI bám theo dòng phát triển của dự án, từ chốt hướng đi đến
triển khai từng phần rồi hoàn thiện báo cáo. Bảng~\ref{tab:prompts_build}
tóm tắt các nhóm prompt ở giai đoạn xây dựng hệ thống, còn
Bảng~\ref{tab:prompts_report} là giai đoạn hoàn thiện báo cáo và slide.

\begin{table}[H]
\centering
\caption{Prompt giai đoạn xây dựng hệ thống}
\label{tab:prompts_build}
\begin{tabularx}{\textwidth}{>{\raggedright\arraybackslash}p{4.4cm} >{\raggedright\arraybackslash}p{3.6cm} X}
\toprule
\textbf{Nhóm prompt} & \textbf{Mục đích} & \textbf{Kết quả} \\
\midrule
Viết và tinh chỉnh spec, kế hoạch pilot & Biến ý tưởng thành việc làm được trong 8--10 ngày & \texttt{PROJECT\_SPEC}, kế hoạch pilot Huế, kế hoạch AI service. \\
Dựng AI service chạy được, sửa import theo monorepo & Khởi động được dịch vụ AI & Sửa import + Docker entrypoint; dịch vụ boot. \\
Định nghĩa knowledge chunk schema & Bảo đảm truy vết nguồn (traceability) & Schema chunk có \texttt{source\_type}, \texttt{review\_status}. \\
Dựng lại retrieval cho Culture Scout & Truy hồi tri thức theo ngữ cảnh & Retrieval theo place/category/language/interests. \\
Hiện thực scoring + giải thích đề xuất (Task 4) & Điểm số tính trong code, LLM chỉ diễn đạt & Scoring tất định + fallback khi mô hình lỗi; test đạt. \\
Crawl nguồn chính thống + UNESCO + ebook & KB có dẫn nguồn, duyệt trước & Registry nguồn vetted + knowledge chunk + validator. \\
Tích hợp weather thật từ OpenWeather theo GPS & Biến weather mock thành ngữ cảnh thật, khóa key ở backend & Weather service chung; storyline/explore dùng weather thật; fallback an toàn khi từ chối GPS. \\
Chốt hạ tầng tự host & Tránh phụ thuộc và hạn mức & Firebase $\to$ PostgreSQL + MinIO + JWT; Gemini $\to$ vLLM tự host. \\
\bottomrule
\end{tabularx}
\end{table}

\begin{table}[H]
\centering
\caption{Prompt giai đoạn hoàn thiện báo cáo và slide (đầy đủ, theo thứ tự)}
\label{tab:prompts_report}
\footnotesize
\begin{tabularx}{\textwidth}{@{}p{0.4cm} >{\raggedright\arraybackslash}p{4.2cm} >{\raggedright\arraybackslash}p{3.3cm} X@{}}
\toprule
\textbf{\#} & \textbf{Prompt} & \textbf{Mục đích} & \textbf{Kết quả} \\
\midrule
1 & Vẽ graph pipeline, chia rõ Frontend/Backend/vLLM layer & Làm rõ kiến trúc tổng thể & Sơ đồ 3 lớp (Hình~\ref{fig:system_pipeline}). \\
2 & ``Bạn tự update nguồn giúp tôi đi'' & Số liệu phải có nguồn thật, không bịa & Xác minh 4 thống kê du lịch, sửa \texttt{ref.bib} trỏ URL thật. \\
3 & ``Commit all changes'' & Lưu thay đổi theo dõi & Commit gỡ report CT cũ. \\
4 & ``Commit features also'' & Theo dõi luôn report & Un-ignore \texttt{docs/report/}, commit source + pdf. \\
5 & ``Push to main'' & Phát hành nhánh chính & Fast-forward + push \texttt{main}. \\
6 & Thêm pseudo-code + JSON mẫu vào §3.3, dựa vào \texttt{data/} & Minh họa schema bằng dữ liệu thật & 6 listing JSON từ dữ liệu thật. \\
7 & Dịch vLLM layer xuống (đang dính Backend) & Sửa bố cục sơ đồ & Tách band vLLM khỏi Backend. \\
8 & Nói rõ hạ tầng: Web$\to$Azure, Backend+vLLM$\to$Thundercompute L40 & Khai báo nền tảng chạy thật & Prose + bảng \texttt{tab:infra}. \\
9 & Làm slide theo style FactChecking, ít chữ & Bộ slide thuyết trình & Deck beamer keyword + diagram. \\
10 & Tải ảnh pain-point (×3 URL) vào slide 2 & Đưa ảnh thật vào slide & 3 ảnh gắn slide 2. \\
11 & Thêm slide terminal vLLM & Minh chứng vLLM chạy & Slide log vLLM. \\
12 & Đổi slide đó sang figure \texttt{vllm\_example.png} & Để tự thay ảnh thật sau & Figure placeholder \texttt{\textbackslash IfFileExists}. \\
13 & Thêm figure homepage §7.1 + slide homepage & Đưa ảnh sản phẩm vào & \texttt{fig:homepage} + slide Web App. \\
14 & Thêm bubble map (overview + detail) vào report \& slide & Minh họa bản đồ thật & \texttt{fig:bubble\_maps} + slide 9. \\
15 & Chỉnh size ảnh bubble map (tràn slide) & Sửa tràn khung & 2 map cạnh nhau, cap height. \\
16 & Thêm \texttt{bevietnam.iamphuckhang.dev} đầu báo cáo & Gắn live demo lên bìa & Hyperlink trang tiêu đề. \\
17 & Gỡ lỗi \texttt{setup.sh} (minio Permission denied / Segfault) & Khắc phục vận hành VM & Chẩn đoán chmod + tải lại binary. \\
18 & Gỡ lỗi KV cache vLLM (0.26 < 0.88 GiB) & Chạy được model trên L40 & Chẩn đoán OOM do VL profiling. \\
19 & Bổ sung pipeline \texttt{vllm\_hosting} sinh knowledge & Trình bày khâu offline & \texttt{fig:offline\_pipeline} + bảng agent + slide. \\
20 & Thêm ảnh mobile/storyline/vLLM vào report \& slide & Minh chứng sản phẩm & \texttt{fig:storyline}, \texttt{fig:mobile}, \texttt{fig:vllm\_serve}. \\
21 & Thêm mục trình bày cách prompt qua các session & Minh bạch cách dùng AI & Chương này + tài liệu nhật ký prompt. \\
\bottomrule
\end{tabularx}
\end{table}

\subsection{Nguyên tắc Prompt Vận hành}
\contrib

Nhóm thống nhất một số nguyên tắc khi đặt câu lệnh, để giữ AI bám đúng
việc và để kết quả luôn kiểm chứng được:

\begin{itemize}
  \item \textbf{Nêu rõ mục tiêu và bối cảnh} trong mỗi prompt, ưu tiên đọc
  mã nguồn trước khi kết luận thay vì đoán mò.
  \item \textbf{Thay đổi nhỏ, bám pattern sẵn có} của repo; thực hiện
  triển khai thật trong workspace thay vì chỉ đề xuất.
  \item \textbf{Không bịa nguồn}: mọi số liệu, trích dẫn phải truy được về
  tài liệu thật, nếu không thì bỏ.
  \item \textbf{Bảo toàn thay đổi của nhóm}: không tự ý revert file người
  dùng, không dùng lệnh phá hủy khi chưa được yêu cầu.
  \item \textbf{Làm theo vòng lặp}: AI đề xuất, nhóm biên dịch/chạy
  thử/đọc lại, rồi prompt tiếp để sửa.
\end{itemize}

Tóm lại, AI được dùng chủ yếu ở vai trò \textit{tăng tốc thực thi} (sinh
mã, dựng sơ đồ, soạn báo cáo, gỡ lỗi) chứ không thay nhóm ra quyết định
thiết kế hay tự ý đưa ra số liệu. Nhóm đọc lại, chỉnh sửa, kiểm thử và
chịu trách nhiệm cuối cùng cho mọi dòng mã cũng như nội dung báo cáo.
